\documentclass[11pt,a4paper,sans]{moderncv}

% Style et couleur de moderncv
\moderncvstyle{classic}
\moderncvcolor{green}


% Encodage et mise en page
\usepackage[utf8]{inputenc}
\usepackage[scale=0.85]{geometry}
\setlength{\hintscolumnwidth}{3.0cm}

% Inclusion du paquet moderntimeline
\usepackage[firstyear=2018,lastyear=2026]{moderntimeline}

\usepackage{pdfpages}
\usepackage{amsmath}
\usepackage{amssymb}

\usepackage{enumitem}

% \usepackage{doi}

\newcommand*{\doi}[1]{\href{http://dx.doi.org/#1}{DOI: #1}}

% Informations personnelles
\firstname{Ari Edward}
\familyname{RAPPAPORT}
\title{
  Curriculum vitae
}
\homepage{aerappa.github.io}

\begin{document}

\makecvtitle

\section{Career Timeline}
\tlcventry{2026}{0}{Postdoctoral Researcher}{Laboratoire Jacques-Louis Lions}{Paris, France}{}{}
\tlcventry{2024}{2026}{Postdoctoral Researcher}{ENSTA Paris}{Palaiseau, France}{}{}
\tlcventry{2024}{2026}{Teaching Assistant (Tutorials)}{ENSTA Paris}{Palaiseau, France}{}{}
\tlcventry[blue]{2021}{2024}{PhD Candidate}{Inria Paris}{Paris, France}{}{}
\tlcventry[blue]{2018}{2020}{Master’s Degree}{University of New Mexico}{Albuquerque, United States}{}{}

\section{Education}
\cventry{2021--2024}{PhD in Applied Mathematics}{Sorbonne University}{Paris, France}{\textit{SERENA Team}}{%
Dissertation: \textit{A posteriori error estimates and adaptivity in numerical approximation of PDEs: regularization, linearization, discretization, and floating-point accuracy} \\
Advisor: Martin Vohralík \\
Co-advisor: François Févotte
}
\cventry{2018--2020}{Master’s Degree in Applied Mathematics}{University of New Mexico}{Albuquerque, United States}{}{%
Thesis: \textit{An a posteriori error analysis of stationary incompressible magnetohydrodynamics} \\
Advisors: Jehanzeb Chaudhry and John Shadid
}
\cventry{2013--2018}{Bachelor’s Degree in Computational Mathematics}{University of New Mexico}{Albuquerque, United States}{}{%
Thesis: \textit{Exploration of Black Box Multigrid for Linear Systems Modelling Resistive MHD} \\
Advisors: Jehanzeb Chaudhry, Raymond Tuminaro, and John Shadid
}

\section{Publications}
\subsection{Published Articles}

\begin{enumerate}
    \item \underline{Rappaport, A. E.}, Chaumont-Frélet, T., and Modave, A.
    A hybridizable discontinuous Galerkin method with transmission variables for time-harmonic electromagnetic problems.
    \textit{SIAM Journal on Scientific Computing} \textbf{48}(4) (2026), B629--B650. \\
    \doi{10.1137/25M1756818}
    \item Harnist, A., Mitra, K., \underline{Rappaport, A.}, and Vohralík, M.
    Robust augmented energy a posteriori estimates for Lipschitz and strongly monotone elliptic problems.
    \textit{IMA Journal of Numerical Analysis}, drag036 (2026). \\
    \doi{10.1093/imanum/drag036}
    \item Févotte, F., \underline{Rappaport, A.}, and Vohralík, M.
    Adaptive regularization for the Richards equation.
    \textit{Computational Geosciences} \textbf{28}(6) (2024), 1371--1380. \\
    \doi{10.1007/s10596-024-10309-7}
    \item Févotte, F., \underline{Rappaport, A.}, and Vohralík, M.
    Adaptive regularization, discretization, and linearization for nonsmooth problems based on primal--dual gap estimators.
    \textit{Computer Methods in Applied Mechanics and Engineering} \textbf{418} (2024), 116558. \\
    \doi{10.1016/j.cma.2023.116558}
    \item J.~H. Chaudhry, \underline{A.~E. Rappaport}, and J.~N. Shadid,
    An a posteriori error analysis for the equations of stationary incompressible magnetohydrodynamics,
    \textit{SIAM Journal on Scientific Computing} \textbf{43} (2021), no.~2, B354--B380. \\
    \doi{10.1137/20M1342975}
\end{enumerate}

\subsection{Submitted Articles and Preprints}

\begin{enumerate}[itemsep=1em]
    \item Bonazzoli, M., Ciarlet Jr., P., Modave, A., and \underline{Rappaport, A. E.}
    Analysis of a two-level domain decomposition preconditioner for the time-harmonic Maxwell equations in anisotropic media.
    HAL Preprint 05572507, 2026.

    \item Févotte, F., \underline{A. Rappaport}.
    A thread-parallel implementation of equilibrated flux reconstruction in Julia.
    HAL Preprint 04537446, 2024.
\end{enumerate}

\section{Research and Teaching Experience}
\cventry{Sep 2026--Present}{Postdoctoral Researcher}{Laboratoire Jacques-Louis Lions, Sorbonne Universit\'e}{Paris, France}{%
Research on phase-space methods within the ERC Synergy project PSINumScat (Phase-space-inspired Numerical Methods for High Frequency Wave Scattering).
}{}
\cventry{Sep 2024--Sep 2026}{Teaching Assistant (Tutorials and Labs)}{ENSTA Paris}{Palaiseau, France}{%
Prepared and led tutorial and laboratory sessions, developed pedagogical materials and problem sets, delivered theoretical summaries, and graded student assessments.
}{}
\cventry{Nov 2017--Oct 2020}{Research Intern}{Sandia National Laboratories}{Albuquerque, United States}{%
Department 1442, Computational Mathematics.
Research on a posteriori error analysis for stationary incompressible magnetohydrodynamics.
Project on using Long Short Term Memory (LSTM) methods to enhance POD-based reduced-order models in MHD.}{}
\cventry{May 2015--Jun 2018}{Research Intern}{Voss Scientific}{Albuquerque, United States}{%
Extension of a numerical viscosity algorithm for a PIC code.
Development of interfaces for multigrid Maxwell solvers and adaptation of a preconditioner for an augmented Maxwell system.}{}
\cventry{Aug 2017--May 2018}{Teaching Assistant (Laboratory Sessions)}{Department of Computer Science, University of New Mexico}{Albuquerque, United States}{%
Designed and led computer laboratory sessions, developed progressive hands-on exercises, and provided individual student support to address technical and conceptual difficulties.
}{}

\section{International Conference Presentations}
\begin{itemize}
  \item \href{https://sites.cwi.nl/en/scs/events/wave-ned-conference-on-analysis-modeling-and-numerics-of-waves/}{WAVE-NED Conference on Analysis, Modeling and Numerics of Waves} (Delft, The Netherlands, 2026).
  \item \href{https://www.tuwien.at/en/mg/asc/event/cmam2026}{International Conference on Computational Methods in Applied Mathematics (CMAM)} (Vienna, Austria, 2026).
  \item \href{https://canum2026.math.cnrs.fr/en/}{CANUM -- Congr\`es National d'Analyse Num\'erique} (Saint-Jacut-de-la-Mer, France, 2026).
  \item \href{https://www.dd29.polimi.it/}{29th International Conference on Domain Decomposition Methods} (Milan, Italy, 2025).
  \item \href{https://ins.uni-bonn.de/group/gedicke/page/cmam-10}{International Conference on Computational Methods in Applied Mathematics} (Bonn, Germany, 2024).
  \item \href{https://icosahom2023.org/}{International Conference on Spectral and High Order Methods} (Seoul, South Korea, 2023).
  \item \href{https://www.siam.org/conferences-events/past-event-archive/gs23/}{SIAM Conference on Mathematical and Computational Issues in the Geosciences} (Bergen, Norway, 2023).
  \item \href{https://www.wccm2022.org/}{World Congress on Computational Mechanics} (Yokohama, Japan, 2022).
  \item \href{https://www.siam.org/conferences-events/past-event-archive/cse19/}{SIAM Conference on Computational Science and Engineering} (Spokane, United States, 2019).
\end{itemize}


\section{Invited Team and Laboratory Seminars}
\begin{itemize}
  \item Seminar of the Laboratoire de Mathématiques d’Orsay (Orsay, France, 2026)
  \item Seminar of the  Institut de Mathématiques de Marseille (Marseille, France, 2026)
  \item Seminar of the Jean Kuntzmann Laboratory (Grenoble, France, 2026)
  \item Seminar of the AGM Laboratory (Cergy, France, 2025)
  \item Joint LJLL / Inria Seminar (Paris, France, 2025)
  \item IMAG Team Seminar (Montpellier, France, 2025)
  \item UNM Applied Mathematics Seminar Series (Albuquerque, United States, 2020)
\end{itemize}

\section{Honors and Awards}
\begin{itemize}
  \item SIAM CSE Student Travel Award (2019).
  \item UNM Departmental Honors: Mathematics (2018).
  \item UNM Van Dyke Computer Science Scholarship (2017).
  \item UNM Leonard Engineering Scholarship (2016).
\end{itemize}

\section{Computing Skills and Programming}
\begin{itemize}
  \item \textbf{Programming Languages:} Julia, Python, C/C++, MATLAB, Java, Fortran 90.
  \item \textbf{Operating Systems:} Linux, macOS.
  \item Advanced proficiency with \texttt{vim}, \texttt{slurm}, \texttt{make}, etc.
\end{itemize}

\section{Languages}
\begin{itemize}
  \item English: Native
  \item French: Fluent
  \item Spanish: Basic knowledge
  \item Chinese: Basic knowledge
\end{itemize}
\end{document}

